% !TeX root = ../papers/nc-p3.tex
\section{二十点が与える恒等式と十五次元の成分}
\label{sec:twenty-conormal-component}
前節の二つの有理点を$r,t$とする。全テンソルのアフィン空間の中で、
二つの隣接切断が共に階数6となる局所閉スキームを$Z$とする。
以下の恒等式は$Z$を定義する方程式の間に、周囲の環上で成立する。

\begin{proposition}\label{prop:twenty-conormal}
$W$とその第二切断$N=W^\tau$が階数6である点を考える。
第一切断の左六関係を消す$([a],[b])\in\mathbb P^3\times\mathbb P^3$が
横断的な二十点であるとする。各点で、$J_b c=c\otimes b$,
$J_a d=a\otimes d$と書く。$J_b^{\mathsf t}N$と$N J_a$が階数3で、
各々の左核・右核を$c,d$が生成するとする。二十個の
$(c\otimes b)\otimes(a\otimes d)$が線型独立なら、第一切断の
階数6軌跡上で第二切断の百個のSchur方程式は局所エタールに八十個で生成される。
さらに二方程式系を合わせた微分の階数が180なら、$Z$はその点で滑らかな76次元である。
\end{proposition}
\begin{proof}
第一切断の階数6軌跡は、六次小行列を可逆にすれば滑らかな156次元である。
六つの双一次式の横断的零点は射影的な完全交叉で、近傍で有限エタールに動く。
番号付け被覆上で二十対を切断として扱い、座標の局所代表を選ぶ。
第一切断の左核条件から、第二切断の階数に関係なく
\[
 J_b^{\mathsf t}N J_a=0
\]
が四行四列の行列として成立する。第二切断の可逆な六次小行列を左上に移して
\[
 N=P Q+E,\quad
 P=\binom{N_{11}}{N_{21}},\quad
 Q=(I_6\ \ N_{11}^{-1}N_{12}),\quad
 E=\begin{pmatrix}0&0\\0&B\end{pmatrix},\qquad
 B=N_{22}-N_{21}N_{11}^{-1}N_{12}
\]
と書く。$X=QJ_a$, $Y=J_b^{\mathsf t}P$は基点で階数3、$YX=0$なので
$\operatorname{im}X=\ker Y$となる。この三平面の補空間$H$を固定する。
近くで$k^4\xrightarrow{X} k^6/H$と$H\xrightarrow{Y}k^4$は階数3のままである。
前者の核を$d'$、後者の左核を$c'$で局所的に生成する。剰余の階数を仮定せずに
$Xd'\in H$かつ$(c')^{\mathsf t}YH=0$となる。従って
\[
 0=(c')^{\mathsf t}J_b^{\mathsf t}N J_a d'
   =(J_b c')^{\mathsf t} E(J_a d')
\]
が周囲の第一切断軌跡上の正確な恒等式になる。基点で$c',d'$は$c,d$である。
基点の第二切断の左右核を非主行・非主列の十座標へ射影する写像は、どちらも
六次主行列の可逆性から同型である。従って$B$の百係数の前に並ぶ二十行は
基点で独立である。対応する二十次小行列を可逆にすれば、二十個の$B$の成分が
残る八十個の線型結合となる。両切断を合わせた微分の階数が180なら、
この八十方程式の第一軌跡上の微分は全射である。Jacobi判定とエタール降下で結論を得る。
\end{proof}

\begin{theorem}\label{thm:full-fifteen-component}
$r,t$を含む$Z$の既約成分は両点で一意で、同一である。その添字保持同型を除いた
次元は15である。幾何学的一般点は標準Koszul左右AS正則、最小周期4、反対階数16で、
四位置の射影自己同型群は自明である。$r$の近くで片側対称窓の基底変換による像は
滑らかな因子をなし、その上の反対階数は10、外側では16である。
反対行列式はこの因子でちょうど六重に消える。
\end{theorem}
\begin{proof}
$r,t$の第一辺は六つの双一次式の被約二十点完全交叉である。
素数12619と2657で四配置を実際にたどり、二つの評価階数が各点で3、
二十行二百五十六列の上記行列が階数20であることを
\path{computations/verify_twenty_conormal_identities.py}で検算する。
第一辺の有限エタールな番号付けに沿って階数下界は標数零に延びる。
基点では$YX=0$なので両方の階数が3以上ならともに3であり、
核直線も特殊化と可換である。前節によれば各点の接空間は$61+15=76$次元で、
微分の階数は180である。命題\ref{prop:twenty-conormal}の仮定が全て成立する。

節\ref{sec:klein-elliptic-completion}の既約族の閉包は$\lambda=1$に
第二補完の一般点を含む。四基底の変換と全体倍率を許した閉包は既約であり、
第二補完の有理点$r$と楕円曲線族の有理点$t$をともに含む。
両点の局所環が正則なので一意な成分が一致する。一般の軌道は、四位置の
点配置が有限で射影枠を含むためアフィン次元61である。よって有理商は
$76-61=15$次元となる。ここでの滑らかさはテンソル軌跡の主張であり、
有限安定群を持つ点での粗商の滑らかさを前提としない。
正則点$t$から有限Frobenius条件と内部次数別の完全性を一般点に移す。
$t$の四配置の非同値性が一般点で保たれるので最小周期は4、反対階数も開条件で16となる。
$t$の射影安定群は四元群である。各置換を枠から復元して微分する前節の議論を、
今度は$Z$の滑らかな近傍で用いる。接空間の軌道商には全非自明文字が現れるため、
一般点での射影安定群は自明である。

最後に偶二位置を同定して対称となる条件を考える。基点$r$で反対切断の像は
対称二乗全体である。階数が10の近傍をこの対称な部分で取ると、同定は射影的に一意である。
これを動かす自由度は$\mathrm{PGL}_4$の15次元である。より具体的には、
対称二乗平面の基底変換軌道はGrassmann多様体の局所閉な斉次空間で、
その安定群は対角$\mathrm{PGL}_4$である。階数10の部分では像平面が正則に定まるので
同定も正則に回復できる。従って局所的に、この片側対称軌跡は滑らかな60次元の
固定同定軌跡に同定の15変数を加えた75次元の滑らかな部分多様体となる。
実際、反対階数が10以上の開集合で階数10は閉条件で、像平面の斉次軌道も
その各点の近くで閉にできる。$Z$は滑らかな76次元なので、$r$の近くでこの因子の
局所方程式を$u$とできる。反対切断の十次小行列でSchur補行列を取ると、
残る六行六列の各成分は因子上で零なので$u$で割れる。前節の横断方向の
核から余核への微分は可逆である。よって$u$で割った六行六列行列は
$r$の近傍で可逆となり、階数の二分と行列式の消滅次数が従う。
\end{proof}
